\begin{ouverture}[calcul]{2}{Additionner, soustraire, multiplier}{Partie 1 · Nombres entiers et calculs}
{Techniques posées · Calcul réfléchi · Parenthèses}
\objectif{Poser et contrôler sommes, différences et produits.}
\objectif{Calculer de tête avec des procédures expliquées.}
\objectif{Respecter les parenthèses dans un calcul.}
\prerequis{Les quatre opérations du CM1.}
\end{ouverture}
\pagestyle{otMain}\essentiel\fontsize{11}{15}\selectfont
\phantomsection\label{cours-2-1}
\essSection{1}{Consolider les techniques}
\begin{savoirfaire}[Poser, estimer, contrôler]Aligner les rangs, gérer les retenues, estimer avant de calculer : $38\,675+4\,850$ est proche de $38\,700+4\,900=43\,600$ ; le calcul posé donne $43\,525$. Une soustraction se contrôle par l’addition : $70\,200-2\,460=67\,740$ car $67\,740+2\,460=70\,200$.\end{savoirfaire}
\begin{exemple}$485\times36$ : $485\times6=2\,910$ et $485\times30=14\,550$ ; la somme donne $17\,460$. L’estimation $500\times35=17\,500$ confirme l’ordre de grandeur.\end{exemple}
\begin{propriete}Multiplier des nombres ronds : $300\times40=12\,000$, car $3\times4=12$ suivi des trois zéros. Cette règle vient des rangs : des centaines multipliées par des dizaines donnent des milliers.\end{propriete}
\begin{savoirfaire}[Des procédures de tête]Ajouter ou retirer 98 ou 99 : passer par 100. $457+99=457+100-1=556$ ; $623-98=623-100+2=525$. Les doubles aident aussi : $36+37=36+36+1=73$.\end{savoirfaire}
\clearpage
\phantomsection\label{cours-2-2}
\essSection{2}{Les parenthèses}
\begin{definition}Les parenthèses s’effectuent d’abord : $5\times(12-4)=5\times8=40$, alors que $5\times12-4=56$. Elles peuvent s’emboîter par paires : $(20-(3+5))\times2=(20-8)\times2=24$.\end{definition}
\begin{exemple}Traduire « le double de la somme de 8 et 7 » : $2\times(8+7)=30$. Sans parenthèses, $2\times8+7=23$ dit autre chose : « le double de 8, puis 7 de plus ».\end{exemple}
\begin{savoirfaire}[Une ligne par étape]Dans un calcul à étapes, écrire une ligne par transformation et recopier ce qui ne change pas : $(35+15)\times(9-5)=50\times4=200$.\end{savoirfaire}
\begin{attention}Une expression se lit en entier avant de calculer. Repérer les paires de parenthèses : chaque ouvrante a sa fermante, et l’intérieur se traite en premier.\end{attention}

\clearpage\rubrique{Découverte}
\begin{blocLivre}{ex-2-1}\begin{exercice}[1]{Deux sommes posées}
Poser et calculer $38\,675+4\,850$ et $256\,908+87\,314$.
\end{exercice}\end{blocLivre}
\begin{blocLivre}{ex-2-2}\begin{exercice}[1]{Deux différences posées}
Poser et calculer $70\,200-2\,460$ et $500\,000-123\,456$.
\end{exercice}\end{blocLivre}
\begin{blocLivre}{ex-2-3}\begin{exercice}[1]{Presque cent}
Calculer de tête $457+99$, $623-98$ et $350+198$.
\end{exercice}\end{blocLivre}
\begin{blocLivre}{ex-2-4}\begin{exercice}[1]{Des nombres ronds}
Calculer $300\times40$, $70\times800$ et $600\times50$.
\end{exercice}\end{blocLivre}
\begin{blocLivre}{ex-2-5}\begin{exercice}[1]{Un produit posé}
Poser et calculer $485\times36$.
\end{exercice}\end{blocLivre}

\rubrique{Entraînement}
\begin{blocLivre}{ex-2-6}\begin{exercice}[2]{Les parenthèses d’abord}
Calculer $5\times(12-4)$, puis $5\times12-4$. Pourquoi les résultats diffèrent-ils ?
\end{exercice}\end{blocLivre}
\begin{blocLivre}{ex-2-7}\begin{exercice}[2]{Deux paires}
Calculer $(35+15)\times(9-5)$, puis $(20-(3+5))\times2$.
\end{exercice}\end{blocLivre}
\begin{blocLivre}{ex-2-8}\begin{exercice}[2]{Traduire une phrase}
Écrire en une seule expression puis calculer : « le triple de la différence entre 40 et 13 ».
\end{exercice}\end{blocLivre}
\begin{blocLivre}{ex-2-9}\begin{exercice}[2]{La bonne expression}
Trois amis achètent chacun un sandwich à 4 euros et une boisson à 2 euros. Parmi $3\times4+2$ et $3\times(4+2)$, choisir l’expression qui donne la dépense totale, puis calculer.
\end{exercice}\end{blocLivre}
\begin{blocLivre}{ex-2-10}\begin{exercice}[2]{Compléter les étapes}
Compléter : $56\times\square=5\,600$ ; $\square-4\,500=12\,800$ ; $7\times\square=4\,900$.
\end{exercice}\end{blocLivre}

\rubrique{Problèmes}
\begin{blocLivre}{ex-2-11}\begin{exercice}[3]{Le festival}
Un festival a vendu $12\,760$ billets le samedi et $9\,845$ le dimanche. La salle pouvait accueillir $25\,000$ personnes au total sur le week-end. Combien de places sont restées libres ?
\end{exercice}\end{blocLivre}
\begin{blocLivre}{ex-2-12}\begin{exercice}[3]{La commande de la mairie}
La mairie commande 36 tables à 485 euros l’une et 144 chaises à 59 euros l’une. Quel est le montant total de la commande ?
\end{exercice}\end{blocLivre}
\begin{blocLivre}{ex-2-13}\begin{exercice}[3]{Une facture à vérifier}
Une facture annonce : « 24 cartouches à 18 euros : 532 euros ». Vérifier par une estimation puis par le calcul exact, et corriger si besoin.
\end{exercice}\end{blocLivre}
\begin{blocLivre}{ex-2-14}\begin{exercice}[3]{Le double piège}
Zoé calcule $2\times15+5=40$ et Tom calcule $2\times(15+5)=40$. L’un des deux se trompe : lequel ? Donner les deux bons résultats.
\end{exercice}\end{blocLivre}
\begin{blocLivre}{ex-2-15}\begin{exercice}[3]{Le tour de magie}
Un magicien annonce : « Choisis un nombre, ajoute 25, retire 17, retire ton nombre de départ : je connais ton résultat ! » Tester avec 40 et avec 358, puis expliquer.
\end{exercice}\end{blocLivre}

\clearpage
